\documentclass[10pt,a4paper]{amsart}
\usepackage[margin=19mm]{geometry}
\usepackage{amsmath,amssymb,mathtools}
\usepackage[scr=rsfs,cal=euler]{mathalfa}
\usepackage{xfrac}
\usepackage[unicode,hidelinks,bookmarks=false]{hyperref}
\newcommand{\ie}{i.e.\ }
\newcommand{\cf}{cf.\ }
\newcommand{\resp}{respectively\ }
\newcommand{\Subsection}{\subsection}
\providecommand{\nobreakdash}{\nobreak\hbox{-}}
\setlength{\emergencystretch}{2em}
\multlinegap0pt
\usepackage{mathtools}
\usepackage{amssymb}
\usepackage{bm}
\usepackage[shortlabels]{enumitem}
\newtheorem{theorem}{Theorem}
\newtheorem{lemma}{Lemma}
\newcommand{\Gr}{\operatorname{Gr}}
\newcommand{\PP}{\mathbb P}
\newcommand{\Pl}{\operatorname{Pl}}
\newcommand{\Span}{\operatorname{span}}
\newcommand{\cU}{\mathcal U}
\title{Grassmann--Pl\"ucker Parametrization of Convolutional Filter Subspaces:\ Regularity and Closed Embeddings}
\author{Hongyu Yuan, Huaiqing Zuo}
\date{Source: arXiv:2609.03361v1 (CC BY 4.0)}
\begin{document}
\maketitle

\noindent\emph{Source and licence.} Grassmann--Pl\"ucker Parametrization of Convolutional Filter Subspaces:\ Regularity and Closed Embeddings. Version arXiv:2609.03361v1 (CC BY 4.0), distributed by arXiv under the Creative Commons Attribution 4.0 International licence, http://creativecommons.org/licenses/by/4.0/, as recorded in the arXiv licence metadata for this exact version and shown on the abstract page https://arxiv.org/abs/2609.03361v1. Full licence notice: \texttt{licenses/arxiv-2609.03361-CC-BY-4.0.txt}.\par\smallskip

\noindent\textit{Editorial scope.} Counted unit: the article's lemma lem:subgrassmannian-closed (source0.tex lines 938--941). The article's complete original proof of it is delivered with it (source0.tex lines 943--1058, one \texttt{proof} environment of 116 lines). No mathematics is written, repaired or re-derived: the statement and the proof are the article's own lines, in the article's own notation, and no retained line is substituted or re-flowed.  Retained uncounted as the source-local context the proof needs: lines 577--585.  Nothing was omitted from any retained span, so the package carries no omission record.  No retained line contains a citation, so the wrapper carries no bibliography and no bibliography entry.  No retained line contains a figure environment, an image-inclusion command or drawing code, so no image is delivered and the package declares no asset dependency.  Editorial formatting only, in addition: an AMS \texttt{amsart} preamble that restates the article's own declarations and macros used by the retained lines, namely the environments \texttt{theorem}, \texttt{lemma} on separate counters, together with the standard packages those lines need.  The retained environments print in this document as Theorem 1, Lemma 1 in order of appearance, whereas the article prints the same items with the numbering of its own full document.  The complete native source of the article as distributed in the arXiv e-print for this exact version is bundled unchanged as \texttt{sources/209290/source0.tex}.
\begin{theorem}[Pl\texorpdfstring{\"u}{u}cker closed embedding]
\label{thm:Plucker-closed}
The Pl\"ucker map
\[
  \Pl_V:\Gr(q,V)\longrightarrow\PP(\wedge^qV)
\]
is a closed embedding. Its image is the projective Grassmann variety, classically
defined by the quadratic Pl\"ucker relations.
\end{theorem}

\begin{lemma}[Closed embedding of a sub-Grassmannian]
\label{lem:subgrassmannian-closed}
The map $j$ is a closed embedding.
\end{lemma}

\begin{proof}
Write
\[
  n=\dim W,\qquad m=\dim H.
\]
Choose a basis $e_1,\ldots,e_n$ of $W$ and extend it to a basis
\[
  e_1,\ldots,e_n,e_{n+1},\ldots,e_m
\]
of $H$.

We first prove that $j$ is injective. If $S_1,S_2\in\Gr(q,W)$ and
$j(S_1)=j(S_2)$, then $S_1$ and $S_2$ are equal as linear subspaces of $H$,
and hence $S_1=S_2$. Thus $j$ is injective.

We next describe its image. Define
\begin{equation}\label{eq:Z-subgrassmannian}
  Z=\left\{
  S\in\Gr(q,H):
  p_J(S)=0\text{ for every }J\nsubseteq\{1,\ldots,n\}
  \right\}.
\end{equation}
Here $J\nsubseteq\{1,\ldots,n\}$ means that $J$ contains at least one index
greater than $n$. We prove that
\begin{equation}\label{eq:image-j-Z}
  j(\Gr(q,W))=Z.
\end{equation}

If $S\in\Gr(q,W)$, every basis vector of $S$ belongs to
$\Span(e_1,\ldots,e_n)$, so
\[
  \wedge^qS\subseteq\wedge^qW.
\]
Therefore all Pl\"ucker coordinates containing an index greater than $n$ vanish,
and hence $j(S)\in Z$. This proves
\[
  j(\Gr(q,W))\subseteq Z.
\]

Conversely, take $S\in Z$. The Pl\"ucker coordinates of $S$ are not all zero,
and every coordinate containing an index greater than $n$ is zero. Hence there
exists
\[
  I=\{i_1<\cdots<i_q\}\subseteq\{1,\ldots,n\}
\]
with $p_I(S)\neq0$. Thus $S$ lies in the standard open set $\cU_I$. By
reordering only $e_1,\ldots,e_n$, we may assume that
$I=\{1,\ldots,q\}$. After normalizing $p_I$ to $1$, the subspace $S$ has a
unique row-space matrix
\begin{equation}\label{eq:subgrassmannian-chart-matrix}
  M_S=
  \begin{pmatrix}
    I_q&A&B
  \end{pmatrix},
\end{equation}
where the columns of $A$ correspond to $e_{q+1},\ldots,e_n$, and the columns of
$B$ correspond to $e_{n+1},\ldots,e_m$.

Let $b_{r\ell}$ be the entry of $B$ in row $r$ and in the column corresponding
to $e_\ell$, where $1\leq r\leq q$ and $n+1\leq\ell\leq m$. Computing the
corresponding minor of \eqref{eq:subgrassmannian-chart-matrix} gives
\[
  p_{\{1,\ldots,\widehat r,\ldots,q,\ell\}}(S)
  =(-1)^{q-r}b_{r\ell}p_{\{1,\ldots,q\}}(S)
  =(-1)^{q-r}b_{r\ell}.
\]
This index set contains $\ell>n$. Since $S\in Z$, the left-hand side is zero,
and therefore
\[
  b_{r\ell}=0
\]
for every $r,\ell$. Thus $B=0$, every row of $M_S$ belongs to $W$, and
$S\subseteq W$. Hence $S\in j(\Gr(q,W))$, proving
\eqref{eq:image-j-Z}.

Under the Pl\"ucker embedding, \eqref{eq:Z-subgrassmannian} can be written as
\begin{equation}\label{eq:intersection-subgrassmannian}
  \Pl_H(Z)
  =\Pl_H(\Gr(q,H))\cap\PP(\wedge^qW).
\end{equation}
The space $\PP(\wedge^qW)$ is the projective linear subspace of
$\PP(\wedge^qH)$ defined by the homogeneous linear equations
\[
  p_J=0,
  \qquad J\nsubseteq\{1,\ldots,n\}.
\]
By Theorem~\ref{thm:Plucker-closed}, $\Pl_H(\Gr(q,H))$ is a closed subvariety.
Therefore \eqref{eq:intersection-subgrassmannian} shows that $Z$ is a closed
subvariety of $\Gr(q,H)$.

Finally, we prove that $j$ identifies $\Gr(q,W)$ isomorphically with $Z$. On any
standard open set with $I\subseteq\{1,\ldots,n\}$, points of $\Gr(q,W)$ are
represented by matrices
\[
  \begin{pmatrix}I_q&A\end{pmatrix},
\]
whereas the corresponding points of $Z$ are represented by
\[
  \begin{pmatrix}I_q&A&0\end{pmatrix}.
\]
Thus, in these affine coordinates, $j$ is
\[
  A\longmapsto(A,0),
\]
and its inverse on the image is
\[
  (A,0)\longmapsto A.
\]
Both maps are given by coordinate polynomials and are therefore regular. These
standard open sets cover $\Gr(q,W)$ and $Z$, so
\[
  j:\Gr(q,W)\xrightarrow{\sim}Z
\]
is an isomorphism. Since $Z$ is closed in $\Gr(q,H)$, the map $j$ is a closed
embedding.
\end{proof}

\end{document}
